\definecolor{links}{HTML}{2A1B81}
\hypersetup{colorlinks,linkcolor=,urlcolor=links}

\usetheme{Boadilla}
\usecolortheme{seahorse}
\usefonttheme{serif}
\beamertemplatenavigationsymbolsempty

\usepackage{luacode}
\usepackage{luatexja}
\usepackage{pgfpages}
\usepackage[osf]{mathpazo}

\begin{luacode*}
  USE_IPAFONT = os.getenv"USE_IPAFONT"
  USE_YUFONT = os.getenv"USE_YUFONT"
  
  if USE_YUFONT == "true" then
    tex.sprint("\\AtBeginDocument{\\usepackage[yu-osx, deluxe, expert]{luatexja-preset}}")
    tex.sprint("\\AtBeginDocument{\\setmainjfont[BoldFont=YuGothic Bold]{YuGothic Medium}}")
  elseif USE_IPAFONT == "true" then
    tex.sprint("\\AtBeginDocument{\\usepackage[ipaex, deluxe, expert]{luatexja-preset}}")
  else
    tex.sprint("\\AtBeginDocument{\\usepackage[hiragino-pro, deluxe, expert]{luatexja-preset}}")
    tex.sprint("\\AtBeginDocument{\\setmainjfont[BoldFont=HiraKakuPro-W6]{HiraKakuPro-W3}}")
  end
\end{luacode*}

\usepackage{epigraph}
\usepackage{etoolbox}
\usepackage{tikz}
\usepackage{framed}
\usepackage{libertine}
\usepackage[final]{listings}
\usepackage{amsmath}
\usepackage{mathtools}
\usepackage{tikz-qtree}

%\setbeameroption{show notes on second screen=right}

%\usetikzlibrary{arrows,automata,shapes,backgrounds}

\setmainfont[Numbers=OldStyle, BoldFont=Palatino Bold]{Palatino}
\setsansfont{CMU Sans Serif}
\setmonofont{CMU Typewriter Text}

\title[\href{https://github.com/y-yu/extensibleException}{Extensible Exception}]{%
  \href{https://github.com/y-yu/extensibleException}{Extensible Exception} \\
  {\normalsize \href{http://kbkz.connpass.com/event/32420/}{kbkz.tech \#10}} \\
}
\author{吉村 優}
\date{July 16, 2016}
\institute[\url{https://twitter.com/\_yyu\_}]{%
  \url{https://twitter.com/\_yyu\_}\\
  \url{http://qiita.com/yyu}\\
  \url{https://github.com/y-yu}\\
}

\newfontfamily\quotefont[Ligatures=TeX]{Linux Libertine O} % selects Libertine as the quote font

\newcommand*\quotesize{60} % if quote size changes, need a way to make shifts relative
% Make commands for the quotes
\newcommand*{\openquote}
   {\tikz[remember picture,overlay,xshift=-4ex,yshift=-2.5ex]
   \node (OQ) {\quotefont\fontsize{\quotesize}{\quotesize}\selectfont``};\kern0pt}

\newcommand*{\closequote}[1]
  {\tikz[remember picture,overlay,xshift=1.5ex,yshift={#1}]
   \node (CQ) {\quotefont\fontsize{\quotesize}{\quotesize}\selectfont''};}

\newcommand*\shadedauthorformat{\emph} % define format for the author argument

% Now a command to allow left, right and centre alignment of the author
\newcommand*\authoralign[1]{%
  \if#1l
    \def\authorfill{}\def\quotefill{\hfill}
  \else
    \if#1r
      \def\authorfill{\hfill}\def\quotefill{}
    \else
      \if#1c
        \gdef\authorfill{\hfill}\def\quotefill{\hfill}
      \else\typeout{Invalid option}
      \fi
    \fi
  \fi}
% wrap everything in its own environment which takes one argument (author) and one optional argument
% specifying the alignment [l, r or c]
%
\newenvironment{shadequote}[2][l]%
{\authoralign{#1}
\ifblank{#2}
   {\def\shadequoteauthor{}\def\yshift{-2ex}\def\quotefill{\hfill}}
   {\def\shadequoteauthor{\par\authorfill\shadedauthorformat{#2}}\def\yshift{2ex}}
\begin{quote}\openquote}
{\shadequoteauthor\quotefill\closequote{\yshift}\end{quote}}

\newfontfamily\listingfont{Menlo}
\definecolor{dkgreen}{rgb}{0,0.6,0}
\definecolor{gray}{rgb}{0.5,0.5,0.5}
\definecolor{mauve}{rgb}{0.58,0,0.82}

\lstset{
  basicstyle=\listingfont,
  frame=single,
  breaklines=true
}

\lstdefinestyle{scala}{
  basicstyle=\listingfont\scriptsize,
  breakatwhitespace=false,
  language=scala,
  captionpos=b,
  commentstyle=\listingfont\scriptsize\color{dkgreen},
  extendedchars=true,
  xleftmargin=2em,
  keepspaces=true,
  keywordstyle=\listingfont\scriptsize\color{blue},
  emphstyle=\listingfont\scriptsize\color{cyan},
  rulecolor=\listingfont\scriptsize\color{black},
  showspaces=false,
  showstringspaces=false,
  showtabs=false,
  stringstyle=\listingfont\scriptsize\color{mauve},
  tabsize=2
}

\lstdefinelanguage{scala}{
  morekeywords={abstract,case,catch,class,def,%
    do,else,extends,false,final,finally,%
    for,if,implicit,import,match,mixin,%
    new,null,object,override,package,%
    private,protected,requires,return,sealed,%
    super,this,throw,trait,true,try,%
    type,val,var,while,with,yield},
  moreemph={Byte,Short,Int,Long,Float,Double,Char,
    String,Boolean,Unit,Null,Nothing,Any,AnyRef},
  otherkeywords={=>,<-,<\%,<:,>:,\#,@},
  sensitive=true,
  morecomment=[l]{//},
  morecomment=[n]{/*}{*/},
  morestring=[b]",
  morestring=[b]',
  morestring=[b]"""
}

\makeatletter
\def\@fnsymbol#1{\ensuremath{\ifcase#1\or *\or \dagger\or \ddagger\or
   \mathsection\or \mathparagraph\or \|\or **\or \dagger\dagger
   \or \ddagger\ddagger \else\@ctrerr\fi}}
\makeatother

\renewcommand{\thefootnote}{\fnsymbol{footnote}}

\makeatletter
\def\Ldescription{%
  \@ifnextchar[{\beamer@testforospec}{\beamer@descdefault\beamer@descriptionwidth\@@Ldescription}%
}

\def\beamer@testforospec[{\@ifnextchar<{\beamer@scandefaultospec[}{\@Ldescription[}}%

\def\beamer@scandefaultospec[#1]{\def\beamer@defaultospec{#1}\Ldescription}

\def\@Ldescription[#1]{%
\setbox\beamer@tempbox=\hbox{\def\insertdescriptionitem{#1}
  \usebeamertemplate**{description item}}%
\beamer@descdefault\wd\beamer@tempbox\@@description%
}%

\def\@@Ldescription{%
  \beamer@descdefault35pt%
  \list
  {}
  {\labelwidth\beamer@descdefault\leftmargin2.8em\let\makelabel\beamer@Ldescriptionitem}%
  \beamer@cramped%
  \raggedright
  \beamer@firstlineitemizeunskip%
}

\def\endLdescription{\ifhmode\unskip\fi\endlist}
\long\def\beamer@Ldescriptionitem#1{%
  \def\insertdescriptionitem{#1}%
  \hspace\labelsep{\parbox[b]{\dimexpr\textwidth-\labelsep\relax}{%
        \usebeamertemplate**{description item}%
    }}}
\makeatother

\usetikzlibrary{shapes.callouts} 

\pgfkeys{%
    /calloutquote/.cd,
    width/.code                   =  {\def\calloutquotewidth{#1}},
    position/.code                =  {\def\calloutquotepos{#1}}, 
    author/.code                  =  {\def\calloutquoteauthor{#1}},
    /calloutquote/.unknown/.code   =  {\let\searchname=\pgfkeyscurrentname
                                 \pgfkeysalso{\searchname/.try=#1,                                
    /tikz/\searchname/.retry=#1},\pgfkeysalso{\searchname/.try=#1,
                                  /pgf/\searchname/.retry=#1}}
                            }  


\newcommand\calloutquote[2][]{%
       \pgfkeys{/calloutquote/.cd,
         width               = 5cm,
         position            = {(0,-1)},
         author              = {}}
  \pgfqkeys{/calloutquote}{#1}                   
  \node [rectangle callout,callout relative pointer={\calloutquotepos},text width=\calloutquotewidth,/calloutquote/.cd,
     #1] (tmpcall) at (0,0) {\hfil#2\hfil};
  \node at (tmpcall.pointer){\calloutquoteauthor};    
}


\newcommand\ballref[1]{%
\tikz \node[circle, shade,ball color=structure.fg,inner sep=0pt,%
  text width=8pt,font=\tiny,align=center] {\color{white}\ref{#1}};
}

%\everymath{\displaystyle}

\begin{document}

\frame{\maketitle}

\section{自己紹介}
\begin{frame}[fragile]
  \frametitle{自己紹介}
  
  \begin{columns}
    \begin{column}{0.4\textwidth}
      \centering
      \begin{figure}
        \includegraphics[width=0.9\textwidth]{img/bird2x.png}
      \end{figure}
    \end{column}
    \begin{column}{0.6\textwidth}
      \begin{itemize}
        \item<2-> \href{https://www.coins.tsukuba.ac.jp/}{筑波大学 情報科学類} 学士（COINS11）
        \item<3-> 現在はScalaを書く仕事に従事
        \item<4-> エラー処理に関する話をします
      \end{itemize}
    \end{column}
  \end{columns}
\end{frame}

\section{エラー値とは？}

\subsection{エラー値の階層構造}
\begin{frame}[fragile]
  \frametitle{エラー値とは？}

  \uncover<2-> {
    \begin{block}{エラー値}
      \begin{itemize}
        \item<3-> エラーであることを表す値
        \item<4-> 階層構造（木構造）になるのが一般的
      \end{itemize}
    \end{block}
  }

  \uncover<5-> {
    \begin{exampleblock}{エラー値の階層構造の例}
      \begin{center}
        \begin{tikzpicture}[scale=.8]
          \Tree[.RootException
                    HttpException 
                    DbException 
                    [.FileException 
                        ReadException
                        WriteException ] ]
        \end{tikzpicture}
      \end{center}
    \end{exampleblock}
  }

  \begin{center}
    \uncover<6->{
      \begin{tikzpicture}
        \calloutquote[width=8cm,position={(-1.2,-0.2)},fill=green!30,rounded corners]{どうやって\textbf{階層構造}を作る？}
      \end{tikzpicture}
    }
  \end{center}
\end{frame}

\subsection{継承を用いた表現}
\begin{frame}[fragile]
  \frametitle{継承を用いた表現}

  \begin{center}
    \begin{tikzpicture}[scale=.8]
      \Tree[.RootException
                HttpException 
                DbException 
                [.FileException 
                    ReadException
                    WriteException ] ]
    \end{tikzpicture}
  \end{center}

  \uncover<2->{
    \lstinputlisting[style=scala]{./src/Exception.scala}
  }

  \begin{center}
    \uncover<3->{
      \begin{tikzpicture}
        \calloutquote[width=6cm,position={(1.2,-0.2)},fill=blue!30,rounded corners]{どうして継承を使うの？}
      \end{tikzpicture}
    }
  \end{center}
\end{frame}

\section{サブタイピングとEither}
\subsection{サブタイプ多相}
\begin{frame}[fragile]
  \frametitle{サブタイプ多相}

  \begin{block}{}
    \begin{shadequote}[r]{\scriptsize 筑波大学 プログラム言語論\cite{subtype}}
      型$A$が型$B$のsubtype（部分型）のとき、型$B$の式を書くべきところに、型$A$の式を書いても良い。
    \end{shadequote}
  \end{block}

  \begin{center}
    \uncover<2->{
      \begin{tikzpicture}[scale=.8]
        \Tree[.RootException
                  HttpException 
                  DbException 
                  [.FileException 
                      ReadException
                      WriteException ] ]
      \end{tikzpicture}
    }
  \end{center}

  \uncover<3->{
    \begin{exampleblock}{例}
      \begin{itemize}
        \item<4-> RootExceptionを書くべきところにHttpExceptionを書く
        \item<5-> RootExceptionを書くべきところにReadExceptionを書く
      \end{itemize}
    \end{exampleblock}
  }
\end{frame}

\subsection{Either}
\begin{frame}[fragile]
  \frametitle{Either}

  \uncover<2->{
    \lstinputlisting[style=scala, firstline=1, lastline=9]{./src/Either.scala}
  }

  \begin{itemize}
    \item<3-> \lstinline|flatMap|の型パラメータで$AA >: A$を取る
    \item<4-> $AA >: A$は$AA$が$A$のスーパータイプであることを表す
  \end{itemize}

  \uncover<5->{
    \begin{exampleblock}{}
      \lstinputlisting[style=scala, firstline=11, lastline=15]{./src/Either.scala}
    \end{exampleblock}
  }
\end{frame}

\begin{frame}[fragile]
  \frametitle{Either}

  \begin{exampleblock}{}
    \lstinputlisting[style=scala, firstline=17]{./src/Either.scala}
  \end{exampleblock}

  \begin{center}
    \uncover<2->{
      \begin{tikzpicture}
        \calloutquote[width=9cm,position={(-1.2,-0.3)},fill=green!30,rounded corners]{これの型は\lstinline|Left[RootException]|になる}
      \end{tikzpicture}

      \begin{tikzpicture}[scale=.8]
        \Tree[.\textcolor{red}{RootException}
                  HttpException 
                  DbException 
                  [.FileException 
                      ReadException
                      WriteException ] ]
      \end{tikzpicture}
    }

    \uncover<3->{
      \begin{tikzpicture}
        \calloutquote[width=10cm,position={(1.2,-0.3)},fill=blue!30,rounded corners]{型の上ではFileExceptionと\textbf{区別}できなくなった！}
      \end{tikzpicture}
    }
  \end{center}
\end{frame}

\section{階層構造の拡張}
\begin{frame}[fragile]
  \frametitle{階層構造の拡張}

  \begin{center}
    \begin{tikzpicture}[scale=.8]
      \Tree[.RootException
                HttpException 
                DbException 
                [.FileException 
                    ReadException
                    WriteException ] ]
    \end{tikzpicture}

    \uncover<2->{
      \begin{tikzpicture}
        \calloutquote[width=7cm,position={(-1.2,-0.3)},fill=green!30,rounded corners]{\textbf{後から}この階層構造を変更できる？}
      \end{tikzpicture}
    }

    \uncover<3-> {
      \begin{tikzpicture}[scale=.8]
        \Tree[.RootException
                  [.\textcolor{red}{InternalServerErrorException}
                      HttpException 
                      DbException ] 
                  [.FileException 
                      ReadException
                      WriteException ] ]
      \end{tikzpicture}
    }
  \end{center}
\end{frame}

\begin{frame}[fragile]
  \frametitle{階層構造の拡張}

  \begin{center}
    \begin{tikzpicture}
      \calloutquote[width=5cm,position={(1.2,-0.3)},fill=red!30,rounded corners]{\huge\textbf{無理}では？}
    \end{tikzpicture}

    \uncover<2->{
      \begin{tikzpicture}
        \calloutquote[width=11cm,position={(-1.2,-0.3)},fill=blue!30,rounded corners]{\huge\textbf{継承}でやるのはよくない？}
      \end{tikzpicture}
    }

    \uncover<3->{
      \begin{tikzpicture}
        \calloutquote[author={\includegraphics[width=1cm]{./img/bird2x.png}},width=8cm,position={(1.2,-0.5)},fill=green!30,rounded corners]{\huge\textbf{型クラス}でやろう！}
      \end{tikzpicture}
    }
  \end{center}
\end{frame}

\begin{frame}[fragile]
  \frametitle{階層構造の拡張}
  
  \begin{enumerate}
    \item<1-> 新しい型クラス\lstinline|:~>|（Transform）を導入
    \item<2-> 新しいエラー値を定義
    \item<3-> 型クラス\lstinline|:~>|のインスタンスを定義
    \item<4-> 型クラス\lstinline|:~>|のインスタンスを使うようにEitherを拡張
  \end{enumerate}
\end{frame}

\section{新しい型クラスとEitherの拡張}
\subsection{新しい型クラスの定義}
\begin{frame}[fragile]
  \frametitle{新しい型クラス\lstinline|:~>|の定義}

  \uncover<2->{
    \begin{block}{変換を表す型クラス\lstinline|:~>|}
      型$A$から型$B$への変換ができることを表す型クラス
      \uncover<3->{
        \lstinputlisting[style=scala]{./src/Transform.scala}
      }
    \end{block}
  }

  \uncover<4->{
    \begin{exampleblock}{例}
      DbExceptionからRootExceptionへの変換ができることを表す
      \lstinputlisting[style=scala]{./src/Instance.scala}
    \end{exampleblock}
  }
\end{frame}

\begin{frame}[fragile]
  \frametitle{新しいエラー値の定義}

  \uncover<2->{
    \lstinputlisting[style=scala, firstline=1, lastline=2]{./src/InternalServerErrorException.scala}
  }
  
  \uncover<3->{
    \begin{block}{継承に基づく階層構造}
      \begin{center}
        \begin{tikzpicture}[scale=.7]
          \tikzset{level distance=2cm}
          \Tree[.RootException
                    \textcolor{red}{InternalServerErrorException}
                    HttpException
                    DbException
                    [.FileException \edge[roof]; {Read\&WriteException} ] ]
        \end{tikzpicture}
      \end{center}
    \end{block}
  }

  \begin{itemize}
    \item<4-> この時点で、HttpExceptionとDbExceptionはInternalServerErrorExceptionと関係がない
  \end{itemize}
\end{frame}

\subsection{インスタンスの定義}
\begin{frame}[fragile]
  \frametitle{型クラス\lstinline|:~>|のインスタンス定義}

  \uncover<2->{
    \lstinputlisting[style=scala, firstline=4, lastline=16]{./src/InternalServerErrorException.scala}
  }

  \uncover<3->{
    次のインスタンスを定義する
    \begin{itemize}
      \item<4-> DbExceptionからInternalServerErrorExceptionへの変換
      \item<5-> HttpExceptionからInternalServerErrorExceptionへの変換
    \end{itemize}
  }
\end{frame}

\begin{frame}[fragile]
  \frametitle{型クラス\lstinline|:~>|のインスタンス定義}

  \begin{block}{継承に基づく階層構造とインスタンス}
    \begin{center}
      \begin{tikzpicture}[scale=.7]
        \tikzset{level distance=2cm}
        \Tree[.RootException
                  \node(InternalServerError){\textcolor{red}{InternalServerErrorException}};
                  \node(Http){HttpException};
                  \node(Db){DbException};
                  [.FileException \edge[roof]; {Read\&WriteException} ] ]
        \draw[->, color=blue] (Http)..controls +(south west:1.5) and +(south east:1.5)..(InternalServerError);
        \draw[->, color=blue] (Db)..controls +(south west:3) and +(south east:3)..(InternalServerError);
      \end{tikzpicture}
    \end{center}
  \end{block}
\end{frame}

\subsection{Eitherの拡張}
\begin{frame}[fragile]
  \frametitle{Eitherの拡張}

  \begin{block}{既存のEither}
    \lstinputlisting[style=scala, firstline=1, lastline=9]{./src/Either.scala}
  \end{block}

  \uncover<2->{
    Pimp my LibraryパターンでEitherを拡張
    \lstinputlisting[style=scala, firstline=1, lastline=3]{./src/ExceptionEither.scala}
  }

  \begin{itemize}
    \item<3-> \lstinline|map|と\lstinline|flatMap|を拡張
    \item<4-> 新しいメソッド\lstinline|as|を導入
  \end{itemize}
\end{frame}

\begin{frame}[fragile]
  \frametitle{\lstinline|map|と\lstinline|flatMap|の拡張}

  \uncover<2->{
    \lstinputlisting[style=scala, firstline=5, lastline=17]{./src/ExceptionEither.scala}
  }

  \begin{itemize}
    \item<3-> 型クラス\lstinline|:~>|のインスタンスをimplicitパラメータで検索
    \item<4-> \lstinline|Left|の場合、型クラス\lstinline|:~>|のインスタンスを用いて変換
    \item<5-> \lstinline|Right|の場合は元のEitherと同様
  \end{itemize}
\end{frame}

\begin{frame}[fragile]
  \frametitle{新しいメソッド\lstinline|as|の導入}

  \uncover<2->{
    \lstinputlisting[style=scala, firstline=19]{./src/ExceptionEither.scala}
  }

  \begin{itemize}
    \item<3-> 型クラス\lstinline|:~>|のインスタンスをimplicitパラメータで検索
    \item<4-> 型クラス\lstinline|:~>|のインスタンスを用いて変換
  \end{itemize}

  \begin{center}
    \uncover<5->{
       \begin{tikzpicture}
         \calloutquote[width=5cm,position={(1.2,-0.2)},fill=green!30,rounded corners]{これの何が便利なの？}
       \end{tikzpicture}
    }

    \uncover<6->{
       \begin{tikzpicture}
         \calloutquote[width=4cm,position={(-1.2,-0.2)},fill=blue!30,rounded corners]{後で使います}
       \end{tikzpicture}
    }
  \end{center}
\end{frame}

\begin{frame}[fragile]
  \frametitle{例}

  \begin{exampleblock}{}
    \lstinputlisting[style=scala, firstline=1, lastline=9]{./src/Main.scala}
  \end{exampleblock}
\end{frame}

\begin{frame}[fragile]
  \frametitle{既存の階層構造との互換性}

  \uncover<2->{
    \begin{exampleblock}{自らへのインスタンスがない}
      \lstinputlisting[style=scala, firstline=11, lastline=16]{./src/Main.scala}
    \end{exampleblock}
  }

  \uncover<3->{
    \begin{exampleblock}{継承関係との互換性がない}
      \lstinputlisting[style=scala, firstline=18, lastline=25]{./src/Main.scala}
    \end{exampleblock}
  }
\end{frame}

\section{自明なインスタンスの導入}
\begin{frame}[fragile]
  \frametitle{自明なインスタンスの導入}

  \uncover<2->{
    \begin{block}{self（自らへのインスタンス）}
      \lstinputlisting[style=scala, firstline=1, lastline=3]{./src/Implicit.scala}
    \end{block}
  }

  \uncover<3->{
    \begin{block}{superclass（スーパータイプへのインスタンス）}
      \lstinputlisting[style=scala, firstline=5, lastline=7]{./src/Implicit.scala}
    \end{block}
  }
\end{frame}

\section{Transitive}
\begin{frame}[fragile]
  \frametitle{Transitive}
  
  \uncover<2->{
    \lstinline|A :~> B|と\lstinline|B :~> C|から\lstinline|A :~> C|というインスタンスを生成
    \begin{center}
      \begin{tikzpicture}
        \tikzset{node distance=2cm}
        \node (A)              {$A$};
        \node (B) [right of=A] {$B$};
        \node (C) [below of=B] {$C$};
        \path[-stealth]
              (A) edge node {} (B)
              (B) edge node {} (C)
              (A) edge[dashed,->] node {} (C); 
      \end{tikzpicture}
    \end{center}
  }
 
  \uncover<3->{
    \begin{block}{}
      \lstinputlisting[style=scala, firstline=9, lastline=12]{./src/Implicit.scala}
    \end{block}
  }
 
  \begin{center}
    \uncover<4->{
       \begin{tikzpicture}
         \calloutquote[width=4cm,position={(-1.2,-0.2)},fill=green!30,rounded corners]{これで動くの？}
       \end{tikzpicture}
    }
  \end{center}

  \uncover<5->{
    \lstinputlisting[basicstyle=\listingfont\scriptsize]{./src/diverging.txt}
  }
\end{frame}

\begin{frame}[fragile]
  \frametitle{ワンステップの変換を表す型クラスの定義}

  \uncover<2->{
    \begin{block}{\lstinline|:->|}
      \lstinputlisting[style=scala]{./src/Transform1.scala}
    \end{block}
  }

  \begin{itemize}
    \item<3-> 型クラス\lstinline|:->|（Transform1）は推移を含まないワンステップの変換を表す
    \item<4-> 型クラス\lstinline|:->|のインターフェースは\lstinline|:~>|と全く同じ
  \end{itemize}
\end{frame}

\begin{frame}[fragile]
  \frametitle{\lstinline|:~>|のインスタンスを\lstinline|:->|に変更}

  \uncover<2->{
    \begin{block}{Before}
      \lstinputlisting[style=scala, firstline=4, lastline=16]{./src/InternalServerErrorException.scala}
    \end{block}
  }
\end{frame}

\begin{frame}[fragile]
  \frametitle{\lstinline|:~>|のインスタンスを\lstinline|:->|に変更}

  \begin{block}{After}
    \lstinputlisting[style=scala, firstline=18, lastline=30]{./src/InternalServerErrorException.scala}
  \end{block}
\end{frame}


\begin{frame}[fragile]
  \frametitle{Transitiveの定義}
 
  \uncover<2->{
    \begin{block}{Transitive}
      \lstinputlisting[style=scala, firstline=14, lastline=17]{./src/Implicit.scala}
    \end{block}
  }

  \begin{itemize}
    \item<3-> ワンステップの変換\lstinline|:->|を使って発散を防止
  \end{itemize}

  %\uncover<4->{
  %  \lstinputlisting[style=scala, fi]
  %}
\end{frame}

\section*{参考文献}
\begin{frame}
  \frametitle{参考文献}

  \nocite{*}
  \bibliographystyle{junsrt_url}
  \bibliography{ref}
\end{frame}

\begin{frame}
  \frametitle{目次}

  \tableofcontents[hideallsubsections]
\end{frame}

\begin{frame}
  \centering
  {\Huge Thank you for listening!}

  \quad \quad

  {\Huge Any question?}
\end{frame}

\end{document}
